\documentclass[11pt]{article}
\usepackage[T1]{fontenc}
\usepackage[margin=1in]{geometry}
\usepackage{microtype,graphicx,booktabs,tabularx,array,amsmath,amssymb,xcolor,url}
\usepackage[authoryear,round]{natbib}
\usepackage[colorlinks=true,linkcolor=black,citecolor=black,urlcolor=blue]{hyperref}
\newcolumntype{Y}{>{\raggedright\arraybackslash}X}
\newcommand{\EviInitial}{41.14}
\newcommand{\EviPeak}{83.52}
\newcommand{\EviMaxPeak}{84.24}
\newcommand{\EviFinal}{73.22}
\newcommand{\EviCider}{119.77}
\newcommand{\EviMedian}{6.55}
\newcommand{\EviTail}{6.72}
\newcommand{\FullCider}{116.83}
\newcommand{\FullInitialReduction}{86.35}
\newcommand{\FullFinalReduction}{79.94}
\newcommand{\FullPeakReduction}{77.47}
\newcommand{\TopInitialReduction}{18.89}
\newcommand{\TopFinalReduction}{11.57}
\newcommand{\TopPeakReduction}{10.03}
\newcommand{\LookFaster}{86.14}
\newcommand{\PrefixFaster}{11.07}
\newcommand{\TopFaster}{12.35}
\newcommand{\VLFaster}{22.13}

\title{Response to Prof. Fox's Review-Derived Experimental Requirements\\
\large Current EviCastKV Storage-Variant Draft, 6 November 2026}
\author{Private working response for the authors; not the anonymous paper}
\date{}
\begin{document}
\maketitle
\begin{abstract}
This response accompanies a revised manuscript focused on the implemented
shared-source storage variant, EviCastKV-S. It identifies what the completed
256-image COCO study establishes, what was addressed by narrowing the
paper, and what remains experimentally unresolved. It does not assert
universal superiority, byte-matched quality preservation, or validation of
the historical six-component architecture.
\end{abstract}

\section{Source of the concerns and scope of the response}
The available record is the supplied numbered experiment-instruction
package and its description of Prof. Fox's 24 October 2026 assessment,
together with the repository's earlier Fox-review report. The original
private assessment PDF was not available in the inspected files. The
concerns below are therefore \emph{paraphrases of the available
review-derived requirements}, not invented quotations, a verbatim review,
or an official conference rebuttal. Each instruction file is hash-bound
in the paper export's provenance ledger.

The original instructions concern exact prompt selection with unchanged
retained KV entries and append-only decoding. The current implementation
is materially different: it quantizes a shared prompt source, restores
protected exact entries, reconstructs K/V through frozen projections, and
losslessly codes its own generated-history state. The revised paper calls
this the storage variant and makes that change visible in the title,
abstract, design, tables and limitations. It does not say that running
this new study completed the original exact-selection experiment matrix.

The legacy identifier \texttt{mosaickv\_exact} is retained only to locate
raw records. EviCastKV and MosaicKV name the same research project;
EviCastKV-S distinguishes the evaluated variant, not another competing
system. Its prompt approximation is not lossless FullKV, and its native
projection corrections are not uncertainty-guided residual repair.

\section{Completed evidence at a glance}
The immutable producer revision is
\nolinkurl{a6f689fcdfae2b9023045aa412c6d90d8ce79012}.
The common checkpoint is \texttt{llava-hf/llava-1.5-7b-hf}, revision
\nolinkurl{b234b804b114d9e37bb655e11cbbb5f5e971b7a9}.
COCO Caption revision is
\nolinkurl{c3213b54e173ae86517d481de17f31439d6f8c72}.
All methods receive the same 256 images, processor, prompt, BF16 model,
eager backend and 128-prediction fixed schedule. The accepted study
contains 1,536 audited memory rows, 1,536 uninstrumented timing rows,
600 repeated timing rows and six canary workers, for 124 completed GPU
workers. There is no failed accepted sample or FullKV fallback.

\begin{table}[ht]
\centering\small\setlength{\tabcolsep}{4pt}
\resizebox{\textwidth}{!}{% Exported from hash-verified measured records; do not hand edit.
\begin{tabular}{lrrrrrr}
\toprule
Method & Initial & \shortstack{Decode-owned peak\\mean/max} & Final & CIDEr $\uparrow$ & p50 $\downarrow$ & p95 $\downarrow$\\
\midrule
FullKV reference & 301.50 & 370.70/370.70 & 365.00 & 116.83 & 5.85 & 5.97\\
\textbf{EviCastKV-S} & 41.14 & 83.52/84.24 & 73.22 & 119.77 & 6.55 & 6.72\\
Prompt Top-K (A) & 50.73 & 92.83/92.95 & 82.80 & 118.33 & 7.50 & 7.66\\
PrefixKV (R+A) & 45.70 & 86.01/86.56 & 77.78 & 103.10 & 7.38 & 7.54\\
VL-Cache (R+A) & 53.31 & 94.39/94.64 & 85.39 & 122.16 & 13.25 & 17.66\\
LOOK-M (R+A) & 50.72 & 92.82/92.83 & 82.80 & 120.96 & 48.86 & 51.59\\
\bottomrule
\end{tabular}
}
\caption{Measured operating points. Memory is MiB; latency is seconds.
R+A means a local reimplementation plus precision/storage adaptation;
A means an adaptation. The measured byte budgets differ.}
\end{table}

EviCastKV-S has the smallest measured initial, final and mean decode-owned
peak among these compressed points, and lower observed request time.
Its CIDEr is above Prompt Top-K's point estimate, but below VL-Cache's
and LOOK-M's adaptations. The adjusted quality intervals do not establish
preservation against those stronger points. The manuscript reports both
positive and negative findings rather than calling every gate passed.

\section{Point-by-point response}
\subsection{Concern: complete request timing, not selector-only timing}
\textbf{Response: addressed within the current workload.}
The revised evaluation uses complete recomputed requests. Media handling,
processor work, prefill, representation construction, CUDA replay setup,
decode, logits export, detokenization and private-state retirement are
inside the synchronized wall timer. Loaded model weights, writing result
files and evaluating CIDEr are outside it. No timed request borrows
another method's prefill or a graph captured in a previous request.

There are five warmups per method per worker, randomized method order
within each image, and a separate six-image study with 20 paired timed
repetitions. Memory-instrumented timings are not substituted for latency.
The phase table calls its combined region
``prefill/construction/retirement'' rather than presenting it as an
isolated compression kernel. Overlapping host/device intervals are not
summed as exclusive costs.

\textbf{Remaining restriction.} This is single-request HF eager on one
checkpoint. It is not a production-serving or cross-backend timing claim.
The 256-image distributions and six-image repeat intervals are different
experiments and remain separately labeled.

\subsection{Concern: physical storage, padding and ownership must be counted}
\textbf{Response: addressed for the stated initial/final and audited
decode-owned tensor boundaries; incomplete for a whole-request peak.}
The revised storage definition deduplicates backing allocations and
charges headers, exact overlays, patches, index metadata, padding and
live reconstruction buffers. The initial value is measured after cache
binding, not merely from a serialized packet. The final value includes
generated-history growth. The decode audit observes native append overlap
and encoded reconstruction leases, with their distinct observer scopes
recorded explicitly.

The current source result is \EviInitial{} MiB initially,
\EviFinal{} MiB finally and \EviPeak{} MiB mean decode-owned peak,
maximum \EviMaxPeak{} MiB. These are not whole-process allocated or
reserved HBM. We explicitly exclude a complete prefill/construction
KV-peak certificate and do not convert a KV percentage into an identical
total-GPU-memory percentage.

\textbf{Remaining restriction.} A full ownership audit spanning prefill,
construction and vendor-internal resources is still needed for the broad
peak claim. The draft says so rather than marking this item universally
resolved.

\subsection{Concern: equal nominal retention is not an equal-byte comparison}
\textbf{Response: the unsupported claim is withdrawn, not experimentally
resolved by this study.}
The frozen source-slot convention does not produce identical encoded
bytes. EviCastKV-S initially uses 41.14 MiB; the compressed adaptations
use approximately 45.70--53.31 MiB, with different final/peak values.
The paper therefore presents separate operating points. It does not
describe these rows as byte-matched, a swept Pareto frontier, or proof
that the selection policy wins under an identical storage constraint.

The largest initial advantage over a compressed point is approximately
22.82\%, not 30\%. Against the closest-memory PrefixKV adaptation the
final and mean peak differences are only 5.85\% and 2.90\%. Thus no
``30\% smaller at equal quality than every baseline'' statement remains.

\textbf{Remaining restriction.} A shared realized-byte experiment and a
separately registered equal-quality memory analysis remain necessary to
establish those claims. Historical shared-byte selection experiments
cannot be silently transferred to this different storage variant.

\subsection{Concern: representation correctness must be separated from quality}
\textbf{Response: reference distinctions are explicit.}
The paper defines native projection of the same decoded prompt source
as the transport reference. It does not compare a lossy source with
FullKV and call an expected modeling difference a transport defect. The
prompt transport envelope is relative tolerance 0.01 and absolute tolerance
$2^{-6}$ under the pinned dtype predicate. Protected entries cannot hide
mismatches in approximate patches. Generated-history corrections have a
bitwise BF16-word contract relative to the variant's own forward passes.

Six method canaries check all 128 exported logits bitwise between memory
instrumentation and the uninstrumented path. All 1,536 memory/timing
pairs also preserve token IDs and storage endpoints. These tests validate
the observer and runtime path; they are not a proof that prompt quantization
preserves unrestricted FullKV answers.

\textbf{Remaining restriction.} The original masked-subset reference and
retention-one tests answer a different exact-selection contract. The
current numerical envelope must not be presented as an unexplained
relaxation of that contract. Broader adapters need independent checks.

\subsection{Concern: first-token boundary, information access and TTFT}
\textbf{Response: boundary A is declared; boundary B is not claimed.}
The current full-prompt prefill produces the first prediction. Its original
cache is then replaced by the constructed representation for subsequent
forwards. The exporter checks identical first predictions across all
methods. Those first tokens do not demonstrate the compressed cache's
quality. Prompt attention used for protection is available from that one
prefill; benchmark captions and generated-answer correctness do not enter
construction.

Historical final-prompt-token-withheld diagnostics have separate numerical
and no-double-feed requirements. Their failures remain separate; this
paper does not declare them fixed. The present API returns completed
requests rather than streamed tokens, so user-visible emission TTFT is
not reported from an internal prefill-logit timestamp.

\subsection{Concern: faithful and efficient published baselines}
\textbf{Response: provenance and limitations are explicit; native fidelity
is not established for these adaptations.}
Every measured local port is labeled reimplementation plus adaptation.
The additional source/history codec, shape policy and captioning workload
are not the original PrefixKV, VL-Cache or LOOK-M setup. We cite their
papers for the policies but do not attribute the current implementation
cost or failure modes to the published algorithms
\citep{wang2025prefixkv,tu2025vlcache,wan2024lookm}.

There is an especially important quality confound: 42/256 PrefixKV-adaptation
captions reach the fixed output limit; no other method does. Every such
caption remains scored and included. The positive adjusted CIDEr
difference against that adaptation cannot be promoted into a statement
that official PrefixKV loses at its own native PPL or ROUGE-L point.

Likewise, approximately 42.05 seconds in LOOK-M's local combined
prefill/construction region is not evidence that the published LOOK-M
algorithm inherently has that cost. Historical official-path reproduction
and byte-matched results use different configurations and are preserved,
not pooled into this table.

\textbf{Remaining restriction.} Strong original-method comparisons need
efficient faithful paths and native reproduction at their own checkpoint,
metric and operating point. The present paper offers an engineering study
of explicitly adapted local paths, not that broader conclusion.

\subsection{Concern: common optimizations should not favor only EviCastKV}
\textbf{Response: common capability, actual engagement and causal limitation
are all reported.}
All compressed methods request the same uniform-layout replay capability.
The implementation does not silently change adaptive quotas or introduce
padding to make them eligible. The result is that EviCastKV-S engages
the fast path on 256/256 requests and every compressed baseline engages
it on 0/256. This is a genuine layout/runtime difference, not an identical
operator-path experiment.

The new text attributes latency to the combined representation, layout
and execution policy. It does not say that saving a small number of KV
bytes alone explains the gain, nor that a serving engine with paged or
fused attention will achieve the same result. Execution replay graphs
are not the semantic evidence graph.

\textbf{Remaining restriction.} Isolating codec benefit requires additional
layout/execution ablations. Current complete-request latency is useful
but cannot identify that causal contribution by itself.

\subsection{Concern: six-component claims need mechanism ablations}
\textbf{Response: the inactive-design performance claim is removed.}
The historical architecture comprises future-query forecasting, a sparse
cross-modal graph, value-aware utility, budgeted submodular selection,
exact/prototype/residual tiers and uncertainty-guided repair. All six are
inactive in the present study. The paper no longer leads with a three-tier
performance claim or uses CUDA graphs/numerical patches as substitutes
for graph evidence/uncertainty repair.

Its measured contribution is a shared-source format, protected exceptions,
lossless coding of its generated history, and an audited reconstruction
runtime. Existing six-component modules or CPU tests do not validate their
task-quality benefit. Consequently the original routing-fraction,
cross-modal-edge and scoring-component ablation requirements are not
marked completed.

\textbf{Remaining restriction.} Those ablations and model-level numerical
controls are still required if the original algorithmic claims return.
This is a scope correction, not a successful six-component experiment.

\subsection{Concern: confirmation must be fresh at ID and media/source level}
\textbf{Response: not resolved; the experiment is development evidence.}
The 256 images are a frozen, inspected cohort. A historical exclusion
inventory is retained, but global underlying-media/source-group freshness
has not been certified. The original MileBench pools and replacement
source-group uncertainties are not repaired by measuring COCO again.

The revised abstract/evaluation/discussion do not call the result untouched
confirmation. The original proposed COCO-plus-MileBench matrix and the
current six-method, 1,536-row study have different purposes and cannot be
equated because both contain many requests.

\textbf{Remaining restriction.} An ID- and media-audited, preregistered
independent cohort remains necessary before confirmatory generalization
or superiority language. More bootstrap draws cannot create freshness.

\subsection{Concern: corpus statistics, multiplicity and failure denominators}
\textbf{Response: implemented and reported within the inspected cohort.}
All supplied lists of five, six or seven captions are retained. The pinned COCO
evaluator/tokenization ecosystem supplies corpus CIDEr. The 10,000-draw
paired bootstrap recomputes corpus reference document frequencies and
TF--IDF reduction for every resampled image multiset. Shard CIDEr scores
are not averaged. Five quality contrasts use familywise-adjusted intervals.

\begin{table}[ht]
\centering\small\setlength{\tabcolsep}{3pt}
% Exported from hash-verified measured records; do not hand edit.
\begin{tabular}{lrrr}
\toprule
Reference point & $\Delta$CIDEr & Paired 95\% interval & Familywise 95\% interval\\
\midrule
FullKV reference & +2.94 & [+0.005, +6.124] & [-0.915, +7.078]\\
Prompt Top-K (A) & +1.44 & [-1.726, +4.749] & [-2.713, +5.712]\\
PrefixKV (R+A) & +16.67 & [+7.429, +24.811] & [+4.677, +27.616]\\
VL-Cache (R+A) & -2.38 & [-6.386, +1.424] & [-7.805, +2.533]\\
LOOK-M (R+A) & -1.18 & [-5.407, +2.913] & [-6.670, +4.182]\\
\bottomrule
\end{tabular}

\caption{All registered quality contrasts, including negative and
inconclusive results.}
\end{table}

Only the PrefixKV-adaptation contrast has a positive familywise lower
bound. EviCastKV-S versus Prompt Top-K, VL-Cache and LOOK-M does not
establish zero-margin preservation. These are not equivalence tests, and
confidence intervals crossing zero are not called ties or passes.

The six-image repeated study separately reports nested paired timing
ratios and four-contrast adjusted intervals. Its lower time does not
override a failed quality inference. All raw failures from earlier
attempts remain preserved, and accepted denominators are checked rather
than inferred from scheduler exit status alone.

\subsection{Concern: batching, serving throughput and full reproducibility}
\textbf{Response: the artifact is bound; unmeasured scope is excluded.}
The paper exporter verifies producer SHA, immutable ledger, configuration,
raw/manifests, row pairing and generated-table hashes. It independently
recomputes storage endpoints and timing quantiles. No new GPU inference
or manual editing of raw numbers occurs during paper construction.

Batch one is the only model workload. Six concurrent single-GPU jobs
are not batch six, serving throughput or a shared-prefix system. The
optional B=2/4 capacity study, multi-model validation, vLLM/SGLang native
transport and production kernel conclusions remain unmeasured by this
experiment and are not claimed.

\section{Claim changes made in the revised paper}
\begin{enumerate}
\item Replace the historical full six-component performance narrative with
      an explicitly labeled shared-source storage-variant design.
\item Report all six operating points with measured initial, final and
      audited decode-owned peak storage, not nominal ratios alone.
\item Remove same-budget, universal Pareto dominance, 30\% equal-quality
      compression and whole-request peak-certification statements.
\item Retain positive storage and measured local-runtime results, while
      reporting higher VL-Cache/LOOK-M CIDEr and nonconfirmatory intervals.
\item Disclose actual replay engagement, PrefixKV limit-hit captions,
      local baseline adaptations, inspected-cohort status and batch-one
      eager-only hardware/software scope.
\item Keep numerical patches, source quantization, exact overlays and
      lossless generated-history coding conceptually distinct.
\end{enumerate}

\section{What would still be needed for a stronger paper claim?}
The current evidence supports a bounded storage/runtime engineering claim,
not superiority over the strongest baseline. A stronger submission would
need a faithful efficient baseline set at common realized bytes; a complete
whole-request ownership audit; independently audited confirmation; and
quality/memory curves under frozen operating-point rules. Six-component,
multi-model and serving claims additionally need their corresponding
active-path tests and causal ablations. No amount of rewording changes
those empirical requirements.

This response is intentionally candid: some concerns are addressed by new
measurement, some by narrowing the claim, and some remain unresolved.
The authors must review both the scientific interpretation and the
AI-assistance disclosure before using the draft in a submission or
circulating this response as their final position.

\bibliographystyle{plainnat}
\bibliography{references}
\end{document}
